% =============================================================================
% 第6章：诊断——塌缩的深层机制
% =============================================================================

\section{诊断：塌缩的深层机制}
\label{sec:diagnosis}

第\ref{sec:results}节的实验结果确立了四个实证事实：（1）Softmax注意力塌缩为均匀分布；（2）8个注意力头全部同质化；（3）Sparsemax和LOR-MIM分别从不同路径改善了塌缩；（4）注意力多样性与分类性能之间不存在简单的线性关系。本章从嵌入空间结构、损失函数效应和"症状vs病因"三个维度，探究这些事实背后的深层机制。

\subsection{嵌入空间的结构证据}
\label{subsec:embedding_svd}

塌缩的直接原因是Query-Key内积对所有特征对产生近乎相同的值。但其深层原因在更上游——为什么Query和Key向量本身缺乏区分度？

对FTT+Softmax模型训练后的Feature Tokenizer嵌入权重矩阵$\mathbf{W} \in \mathbb{R}^{29 \times 64}$进行奇异值分解（SVD）：第一个奇异值主导了总方差的73.1\%，前三个奇异值合占88.6\%。这意味着29个特征专属的64维嵌入向量几乎全部位于由仅仅3--4个主方向张成的低维子空间中——嵌入层的有效秩远低于其名义维度。

进一步分析任意两个特征嵌入向量$\mathbf{W}_i$与$\mathbf{W}_j$的内积分布：所有$29 \times 28 / 2 = 406$对特征嵌入之间的内积均偏正，且方差极小——没有任何一对特征的嵌入指向相反或正交的方向。在正常的语义嵌入空间中（如词嵌入），"语义相反"或"语义无关"的词对应嵌入向量应呈负内积或近零内积。但在PCA特征空间中，所有嵌入向量指向同一象限——Feature Tokenizer未能通过训练学习到任何"对抗性"的方向差异。

这一嵌入同构性并非特定于本数据集的偶然现象，而是PCA匿名化特征空间的系统性后果：所有V1--V28是同一组主成分的不同线性组合，它们的数值分布相似（标准化后的标准正态），它们的"身份"仅由在主成分空间中的坐标位置区分。Feature Tokenizer面临的任务——从几乎相同的输入分布中学习出截然不同的嵌入方向——在梯度信号层面缺乏足够的区分性压力：所有29个特征对损失函数的梯度贡献是近乎对称的。

置换重要性分析（第\ref{sec:results}节图\ref{fig:permutation_importance}）从另一角度佐证了这一诊断。Softmax、Sparsemax和LOR-MIM+Softmax三个变体的特征排序高度一致（两两Pearson $r > 0.9$），V14虽为首要特征，但其置换效应占比均未超过15\%——即使在特征被随机打乱后，模型性能的下降分散在多个特征上，而非集中于少数关键特征。这种分布式的特征依赖性模式与注意力塌缩的诊断一致：当所有特征嵌入方向趋同时，模型无法形成对单个特征的集中依赖。

综合嵌入SVD分析（第\ref{subsec:embedding_svd}节），塌缩的根因判定为PCA空间中特征嵌入的方向同构性。所有V1--V28是同一组主成分的不同线性组合，Feature Tokenizer从几乎相同的输入分布中学习嵌入方向，梯度信号缺乏区分性压力。置换重要性分析中三个变体高度一致的分布式特征依赖模式（第\ref{sec:results}节图\ref{fig:permutation_importance}）进一步佐证了这一诊断。

\subsection{Sparsemax修复的精确定界：症状而非病因}

综合上述诊断，Sparsemax的+2.6pp AUPRC改善（0.9126$\to$0.9384）可被分解为两个独立效应的叠加：

\begin{enumerate}
    \item \textbf{归一化改善（症状修复）}：Softmax的$e^{z_i} > 0$恒成立特性使其在任何输入上都输出全非零分布——当Query-Key内积高度集中时，输出被迫趋近均匀。Sparsemax将此替换为欧氏投影的截断规则：低于阈值$\tau$的特征被直接赋予零权重。这一替换的效果是解放性的——它让注意力机制摆脱了Softmax在低差异输入上的"强制均匀化"锁死。AUPRC因此从0.9126跃升至0.9384。
    \item \textbf{噪声阻断（校准改善）}：Softmax的微量非零权重在多层的残差累积中被放大——即使单层注意力中1\%的差异权重，在3层Transformer的残差连接中通过逐层累积可能放大至显著水平——最终传导至分类头产生过度自信的预测（ECE=0.2328）。Sparsemax的零截断阻断了这一"噪声累积$\to$过度自信"的传导链（ECE降至0.0120）。
\end{enumerate}

然而，Sparsemax修复的是塌缩的\textbf{症状}（Softmax饱和导致的强制均匀化），而非\textbf{病因}（Feature Tokenizer的嵌入方向缺乏鉴别力）。它允许注意力机制"自由地"在不同特征上分配权重，但\textbf{没有提供关于哪些特征应该被分配更多权重的信息}。如果Feature Tokenizer的嵌入空间中所有特征确实难以区分，Sparsemax的"选择"将退化为基于微小随机差异的任意截断，而非基于语义相关的精准筛选。

这一"症状-病因"的区分解释了为什么Sparsemax后的注意力分布虽然有所改善（不再所有特征均分），但仍然缺乏精准的鉴别性：\textbf{Sparsemax修复了注意力权重的表达自由度，但无法替代Feature Tokenizer为嵌入空间注入区分的语义差异}。

\subsection{LOR-MIM修复的精确定界：数据端对称破缺及其与Sparsemax的拮抗}

LOR-MIM代表第二条修复路径——从数据端打破PCA的全局对称。其原理已在第\ref{sec:background}节详细介绍，此处聚焦于实验结果对诊断框架的意义。

LOR-MIM+Softmax的结果（AUPRC=0.9242, AttnStd=0.01401）确立了第一条结论：\textbf{打破PCA对称确实能让注意力"活过来"——注意力多样性在所有方法中最高（$\times$11.2）}。但AUPRC改善（+1.2pp）远小于Sparsemax（+2.6pp），说明病因端的修复并不能自动转化为分类性能的改善。

\subsubsection{与Sparsemax联合使用的拮抗效应}

更具诊断价值的是LOR-MIM+Sparsemax联合实验（AUPRC=0.9267, AttnStd=0.02064）：注意力多样性被推至$\times$16.5，但AUPRC\textbf{反低于}Sparsemax单独使用（0.9384）。

这一反直觉的结果揭示了两条修复路径之间的\textbf{结构性拮抗}：
\begin{enumerate}
    \item \textbf{Sparsemax的机制依赖logits的序关系}：Sparsemax通过截断低于阈值$\tau$的logits来实现稀疏选择——它选择的是Query-Key内积在原始特征空间中"相对更大"的特征。这一序关系即使在全塌缩的均匀分布中也微弱存在（否则Sparsemax将完全随机选择）。
    \item \textbf{LOR-MIM的混合打乱了这一序关系}：组内线性混合使特征的logits值发生系统性重排——原来"稍高"的特征可能在混合后变成"稍低"，反之亦然。对于Sparsemax而言，这相当于将微弱的信号替换为噪声。
    \item \textbf{拮抗的实质}：LOR-MIM为注意力提供了更多"可区分的差异"，但这些差异的排序结构是随机的——Sparsemax据此做出的稀疏选择因此偏离了数据中的真实判别性结构，导致分类性能反降。
\end{enumerate}

\subsubsection{两条修复路径的互补定位与边界}

综合四模型结果，两条修复路径的精确边界可被界定为：
\begin{itemize}
    \item \textbf{归一化端（Sparsemax）}：解锁注意力权重自由度，依赖数据中残余的微弱序关系进行稀疏选择。优势在于$\times$7.5多样性+AUPRC 0.9384+ECE 0.0120；弱点在于不能提升嵌入层的方向鉴别力。
    \item \textbf{数据端（LOR-MIM）}：在特征层面打破PCA对称，制造可区分的token交互。优势在于多样性恢复最强（$\times$11.2--16.5），且可与任何模型架构组合；弱点在于引入的结构不保证与分类相关，且与Sparsemax稀疏选择机制拮抗。
\end{itemize}

这一分析揭示了论文最核心的诊断发现：\textbf{在PCA匿名化空间中，注意力多样性与分类性能是两个相关但不线性的维度；两条修复路径各自提升其中一个维度，但路径之间存在结构性拮抗，联合使用反而可能削弱整体效果。}

\subsection{塌缩的诊断框架}

基于上述诊断，本文归纳了一个描述性的三阶段诊断模板，用于在其他数据集上检验自注意力是否发生了类似塌缩：

\begin{enumerate}
    \item \textbf{第一阶段（注意力分布检验）}：检查注意力权重分布的标准化标准差$\sigma / (1/d)$——若该值$< 0.05$（即标准差不足均匀基线的5\%），则塌缩高度可能。
    \item \textbf{第二阶段（嵌入结构检验）}：对Feature Tokenizer的嵌入权重矩阵进行SVD——若前3个奇异值占比$> 80\%$，则嵌入空间有效秩不足。
    \item \textbf{第三阶段（修复路径交互检验）}：若多种修复路径可用，检验其联合使用时的协同或拮抗关系——LOR-MIM与Sparsemax的拮抗效应表明，独立有效的修复路径在联合时可能产生负交互。
\end{enumerate}

需要强调的是，该框架在本文中仅在一个数据集上演示，应被视为\textbf{提出的诊断模板}而非已验证的通用框架。此外，LOR-MIM的实验结果暗示了一个潜在的第四阶段——\textbf{修复路径交互检验}：若多种修复路径同时可用，检验它们之间的协同或拮抗关系，可进一步揭示塌缩的机制结构。LOR-MIM+Sparsemax的拮抗效应表明，前三阶段的"归因—修复"逻辑在路径交叉时可能不成立——这也是未来工作中最值得深入的方向。

从金融机构的实践角度，上述诊断框架给出了一个简洁的操作建议：当数据必须经过匿名化处理时，（1）优先使用树模型或不依赖注意力的MLP架构；（2）若业务需求要求使用Transformer类模型，至少应将Softmax替换为Sparsemax（零成本、一行代码、AUPRC +2.6pp）；（3）将注意力分布标准差$\sigma/(1/d)$纳入模型上线前的例行质量检查——塌缩的注意力意味着模型无法为风控决策提供任何有意义的特征归因。
